\begin{exo}{}
	Un nouveau journal communal d'une ville de \np{6000} habitants a imprimé \np{1000} exemplaires de son premier numéro et les a tous vendus au prix de $2\eur{}$ chacun.

	Une étude de marché montre que si le prix du journal baissait d'un certain pourcentage $\frac{x}{100}$ alors les ventes augmenteraient significativement.

	Le but est d'étudier le chiffre d'affaire potentiel pour les prochains numéros en fonction de $x \in \intff{0}{100}$.

	Le comptable a déterminé que le chiffre d'affaire $C$ du journal est défini pour tout $x \in \intff{0}{100}$ par :
	\[
		C(x) = \np{2000} + 80x - x^2
	\]
	
	\begin{enumerate}
		\item \begin{enumerate}
			      \item Démontrer que pour tout $x \in \intff{0}{100}$ :
			            \[
				            C(x) = (-x + 20)(x - 60) + 3\,200
			            \]

			      \item Résoudre l'inéquation $C(x) \ge 3\,200$, puis interpréter le résultat.
		      \end{enumerate}
		\item \begin{enumerate}
			      \item Démontrer que pour tout $x \in \intff{0}{100}$ :
			            \[
				            C(x) < 1\,100 \,\iff\, (-x - 10)(x - 90) < 0
			            \]

					\item En déduire les solutions de l'inéquation :
			            \[
				            C(x) < 1\,100
			            \]

			      \item Interpréter le résultat.
		      \end{enumerate}
	\end{enumerate}
\end{exo}
